\section{Selecting an Oscillator Is Not Starting It}
\label{sec:clocks-starting}

The clock tree has a trap that costs an afternoon the first time and ten minutes every time after:
choosing a source does not turn it on, and a peripheral clocked by a stopped source behaves like a
peripheral that ignores you.

\subsection{The Two-Step That Is Actually Three}

\begin{minted}{rust}
// 1. Ask for it.
dp.cmu_s.clken0().modify(|_r, w| w.hfrcoem23().set_bit());
// 2. Force it on: without this the oscillator idles when nothing demands it,
//    and "nothing" includes the branch you are about to select.
dp.hfrcoem23_s.ctrl().modify(|_r, w| w.forceen().set_bit());
// 3. Wait. This is the step that hangs if the oscillator never comes up.
while !dp.hfrcoem23_s.status().read().rdy().bit_is_set() { /* spin */ }
\end{minted}

Step 2 is the one people miss. It is also the one that makes step 3 a real wait instead of an
instant pass — and an instant pass is worse than a hang, because it looks like success.

\subsection{The Converter That Hung on FSRCO}

The microphone's converter (IADC0) had two separate problems in this project, and one of them was
purely a clock choice. Clocked from \codeinline{FSRCO}, the converter's disable sequence never
completed:

\begin{minted}{text}
# after writing EN = 0 and waiting for the disabling bit to clear
IADC0_STATUS  = ...DISABLING = 1     (still, after 100 ms)
# the next configuration write then faulted, because the peripheral was mid-synchronisation
\end{minted}

The same code, clocked from \codeinline{HFRCOEM23}, disabled in microseconds and never faulted. No
register value changed; only \codeinline{em01grpaclkctrl}. That is the entire difference between a
working microphone driver and an afternoon of bus faults, and it was found by trying the other
oscillator rather than by reading harder.

\subsection{The Rule That Came Out of It}

Before blaming a peripheral, confirm that it is \emph{clocked}, and confirm it in the way the
hardware reports it:

\begin{enumerate}
    \item read the gate bit back — not the value you wrote, the value the register holds;
    \item read the oscillator's ready status, not just its enable bit;
    \item if the peripheral has a synchronisation status, read it: \codeinline{DISABLING},
          \codeinline{SYNCBUSY} and \codeinline{RESETTING} are the hardware telling you the
          difference between ``configured'' and ``ready''.
\end{enumerate}

A peripheral that reads as zero after a correct-looking write has failed one of those three, and the
three have different fixes. Without the read-backs they are indistinguishable — which is why every
driver in this repository returns an error instead of a value when a readiness wait expires.
